% Part: incompleteness
% Chapter: introduction
% Section: definitions

\documentclass[../../../include/open-logic-section]{subfiles}

\begin{document}

\olfileid[pt]{inc}{int}{def}

\olsection{Definições}

Para realizar o projeto de Hilbert de formalizar a matemática e mostrar
que tal formalização é consistente e completa, a primeira providência
seria escolher uma linguagem, um arcabouço lógico e um sistema de
axiomas. Para nossos fins, suponhamos que a matemática possa ser
formalizada em uma linguagem de primeira ordem, isto é, que exista
algum conjunto de !!{constant}s, !!{function}s e !!{predicate}s que,
juntamente com os conectivos e quantificadores da lógica de primeira
ordem, nos permita expressar as afirmações da matemática. A maioria das
pessoas concorda que tal linguagem existe: a linguagem da teoria dos
conjuntos, na qual $\in$ é o único símbolo não lógico. À primeira
vista, é claro, parece muito pouco plausível que uma linguagem tão
simples seja tão expressiva, e foi necessário muito trabalho para
estabelecer que praticamente toda a matemática pode ser expressa nesse
vocabulário tão austero. Para simplificar, por enquanto restrinjamos a
discussão à aritmética, isto é, à parte da matemática que trata apenas
dos números naturais~$\Nat$. A linguagem natural para expressar fatos
da aritmética é~$\Lang L_A$. $\Lang L_A$ contém um único
!!{predicate} de dois lugares~$<$, uma única !!{constant}~$\Obj 0$,
uma !!{function} de um lugar~$\prime$ e duas !!{function}s de dois
lugares,~$+$ e~$\times$.

\begin{defn}
Um conjunto de !!{sentence}s~$\Gamma$ é uma \emph{teoria} se for
fechado sob acarretamento, isto é, se $\Gamma = \Setabs{!A}{\Gamma
\Entails !A}$.
\end{defn}

Há duas maneiras simples de especificar teorias. Uma delas é como o
conjunto de !!{sentence}s verdadeiras em alguma !!{structure}.
Considere, por exemplo, a !!{structure} para $\Lang L_A$ cujo
!!{domain} é~$\Nat$ e na qual todos os símbolos não lógicos são
interpretados da maneira esperada.

\begin{defn}\ollabel{def:standard-model}
O \emph{modelo padrão da aritmética} é a !!{structure}~$\Struct N$
definida da seguinte maneira:
\begin{enumerate}
\item $\Domain N = \Nat$
\item $\Assign{\Obj 0}{N} = 0$
\item $\Assign{\Obj \prime}{N}(n) = n + 1$ para todo $n \in \Nat$
\item $\Assign{\Obj +}{N}(n, m) = n + m$ para todos $n, m \in \Nat$
\item $\Assign{\Obj \times}{N}(n, m) = n\cdot m$ para todos $n, m \in \Nat$
\item $\Assign{\Obj <}{N} = \Setabs{\tuple{n, m}}{n \in \Nat, m \in
  \Nat, n < m}$
\end{enumerate}
\end{defn}

Observe a diferença entre `$\times$' e `$\cdot$': $\times$ é um
símbolo da linguagem da aritmética. Naturalmente, nós o escolhemos para
lembrar a multiplicação, mas $\times$ não é a operação de multiplicação;
é um símbolo de função de dois lugares (oficialmente, $\Obj f^2_1$).
Em contraste, $\cdot$~\emph{é} a função usual de multiplicação. Quando
você vê algo como $n \cdot m$, queremos dizer o produto dos números
$n$ e~$m$; quando vê algo como $x \times y$, estamos falando de um termo
na linguagem da aritmética. No modelo padrão, o símbolo de
função~$\times$ é interpretado como a função $\cdot$ sobre os números
naturais. Para a adição, usamos $+$ tanto como símbolo de função da
linguagem da aritmética quanto como função de adição sobre os números
naturais. Aqui, é preciso usar o contexto para determinar o sentido.

\begin{defn}
A teoria da \emph{aritmética verdadeira} é o conjunto de
!!{sentence}s satisfeitas no modelo padrão da aritmética, isto é,
\[
\Th{TA} = \Setabs{!A}{\Sat{N}{!A}}.
\]
\end{defn}

$\Th{TA}$ é uma teoria, pois, sempre que $\Th{TA} \Entails !A$,
$!A$ é satisfeita em toda !!{structure} que satisfaz~$\Th{TA}$. Como
$\Sat{N}{\Th{TA}}$, temos~$\Sat{N}{!A}$ e, portanto,
$!A \in \Th{TA}$.

A outra maneira de especificar uma teoria~$\Gamma$ é como o conjunto
de !!{sentence}s acarretadas por algum conjunto de sentenças~$\Gamma_0$.
Nesse caso, $\Gamma$ é o ``fecho'' de $\Gamma_0$ sob acarretamento.
Especificar uma teoria dessa forma só é interessante se $\Gamma_0$ for
explicitamente especificado, por exemplo, se os !!{element}s
de~$\Gamma_0$ forem listados. No mínimo, $\Gamma_0$ precisa ser
decidível, isto é, deve existir um teste computável para determinar se
!!a{sentence} conta ou não como elemento de~$\Gamma_0$. Chamamos as
!!{sentence}s de~$\Gamma_0$ de \emph{axiomas} de~$\Gamma$ e dizemos que
$\Gamma$ é \emph{axiomatizada} por~$\Gamma_0$.

\begin{defn}
Uma teoria~$\Gamma$ é \emph{axiomatizada} por~$\Gamma_0$ sse
\[
\Gamma = \Setabs{!A}{\Gamma_0 \Entails !A}
\]
\end{defn}

\begin{defn}
A teoria $\Th{Q}$ axiomatizada pelas sentenças a seguir é conhecida
como ``$\Th{Q}$ de Robinson'' e é uma teoria da aritmética muito
simples.
\begin{align*}
& \lforall[x][\lforall[y][(\eq[x'][y'] \lif \eq[x][y])]] \tag{$!Q_1$}\\
& \lforall[x][\eq/[\Obj 0][x']] \tag{$!Q_2$}\\
& \lforall[x][(\eq[x][\Obj 0] \lor \lexists[y][\eq[x][y']])] \tag{$!Q_3$}\\
& \lforall[x][\eq[(x + \Obj 0)][x]] \tag{$!Q_4$}\\
& \lforall[x][\lforall[y][\eq[(x + y')][(x + y)']]] \tag{$!Q_5$}\\
& \lforall[x][\eq[(x \times \Obj 0)][\Obj 0]] \tag{$!Q_6$}\\
& \lforall[x][\lforall[y][\eq[(x \times y')][((x \times y) + x)]]] \tag{$!Q_7$}\\
& \lforall[x][\lforall[y][(x < y \liff \lexists[z][\eq[(z' + x)][y]])]] \tag{$!Q_8$}
\end{align*}
O conjunto de !!{sentence}s $\{!Q_1, \dots, !Q_8\}$ constitui os
axiomas de $\Th{Q}$; portanto, $\Th{Q}$ consiste em todas as
!!{sentence}s acarretadas por eles:
\[
\Th{Q} = \Setabs{!A}{\{!Q_1, \dots, !Q_8\} \Entails !A}.
\]
\end{defn}

\begin{defn}
Suponha que $!A(x)$ seja !!a{formula} em $\Lang L_A$ com variáveis
livres~$x$ e $y_1$, \dots, $y_n$. Então qualquer !!{sentence} da forma
\[
\lforall[y_1][\dots\lforall[y_n][((!A(\Obj 0) \land \lforall[x][(!A(x)
\lif !A(x'))]) \lif \lforall[x][!A(x)])]]
\]
é uma instância do \emph{esquema de indução}.

A \emph{aritmética de Peano}~$\Th{PA}$ é a teoria axiomatizada pelos
axiomas de $\Th{Q}$ juntamente com todas as instâncias do esquema de
indução.
\end{defn}

\begin{explain}
Toda instância do esquema de indução é verdadeira em~$\Struct{N}$.
Isso é mais fácil de ver quando a !!{formula}~$!A$ tem apenas uma
!!{variable} livre~$x$. Nesse caso, $!A(x)$ define um subconjunto
$X_{!A}$ de~$\Nat$ em~$\Struct{N}$. $X_{!A}$ é o conjunto de todos os
$n \in \Nat$ tais que $\Sat{N}{!A(x)}[s]$ quando $s(x) = n$. A
instância correspondente do esquema de indução é
\[
((!A(\Obj 0) \land \lforall[x][(!A(x) \lif !A(x'))]) \lif 
  \lforall[x][!A(x)]).
\]
Se seu antecedente for verdadeiro em~$\Struct{N}$, então
$0 \in X_{!A}$ e, sempre que $n \in X_{!A}$, o mesmo vale para $n+1$.
Como $0 \in X_{!A}$, obtemos $1 \in X_{!A}$. De $1 \in X_{!A}$,
obtemos $2 \in X_{!A}$, e assim por diante. Logo, para todo
$n \in \Nat$, $n \in X_{!A}$. Mas isso significa que
$\lforall[x][!A(x)]$ é satisfeita em~$\Struct{N}$.
\end{explain}

Tanto $\Th{Q}$ quanto $\Th{PA}$ são teorias axiomatizadas. A grande
questão é: quão fortes elas são? Por exemplo, $\Th{PA}$ pode provar
todas as verdades sobre~$\Nat$ que podem ser expressas em~$\Lang L_A$?
Mais especificamente, os axiomas de $\Th{PA}$ decidem todas as questões
que podem ser formuladas em~$\Lang L_A$?

Outra maneira de formular isso é perguntar: $\Th{PA} = \Th{TA}$?
Evidentemente, $\Th{TA}$ prova (isto é, inclui) todas as verdades
sobre~$\Nat$ e decide todas as questões que podem ser formuladas
em~$\Lang L_A$, pois, se $!A$ é !!a{sentence} em $\Lang L_A$, então ou
$\Sat{N}{!A}$ ou $\Sat{N}{\lnot !A}$ e, portanto, ou
$\Th{TA} \Entails !A$ ou $\Th{TA} \Entails \lnot !A$. Chamamos tal
teoria de \emph{!!{complete}}.

\begin{defn}
Uma teoria $\Gamma$ é \emph{!!{complete}} sse, para toda
!!{sentence}~$!A$ de sua linguagem, ou $\Gamma \Entails !A$ ou
$\Gamma \Entails \lnot !A$.
\end{defn}

\begin{explain}
Pelo Teorema da Completude, $\Gamma \Entails !A$ sse $\Gamma \Proves
!A$; assim, $\Gamma$ é completa sse, para toda !!{sentence}~$!A$ de
sua linguagem, ou $\Gamma \Proves !A$ ou $\Gamma \Proves \lnot !A$.
\end{explain}

Outra questão que somos levados a formular é esta: existe um
procedimento computacional que possamos usar para testar se
!!a{sentence} pertence a~$\Th{TA}$, a $\Th{PA}$ ou mesmo apenas
a~$\Th{Q}$? Podemos tornar isso mais preciso definindo quando um
conjunto (por exemplo, um conjunto de !!{sentence}s) é !!{decidable}.

\begin{defn}
Um conjunto $X$ é \emph{!!{decidable}} sse existe um procedimento
computacional que, com entrada~$x$, devolve~$1$ se $x \in X$ e~$0$ caso
contrário.
\end{defn}

Assim, nossa questão passa a ser: $\Th{TA}$ ($\Th{PA}$, $\Th{Q}$) é
!!{decidable}?

A resposta a todas essas questões será: não. Nenhuma dessas teorias é
decidível. Contudo, esse fenômeno não é específico dessas teorias. Na
verdade, \emph{qualquer} teoria que satisfaça certas condições está
sujeita aos mesmos resultados. Uma dessas condições, satisfeita por
$\Th{Q}$ e $\Th{PA}$, é ser axiomatizada por !!a{decidable} conjunto
de axiomas.

\begin{defn}
Uma teoria é \emph{!!{axiomatizable}} se for axiomatizada por
!!a{decidable} conjunto de axiomas.
\end{defn}

\begin{ex}
Toda teoria axiomatizada por um conjunto finito de !!{sentence}s é
!!{axiomatizable}, pois todo conjunto finito é !!{decidable}. Assim,
$\Th{Q}$, por exemplo, é !!{axiomatizable}.

Teorias axiomatizadas por esquemas, como~$\Th{PA}$, também são
!!{axiomatizable}. Para testar se $!B$ está entre os axiomas
de~$\Th{PA}$, isto é, para computar a função $\Char{X}$ tal que
$\Char{X}(!B) = 1$ se $!B$ é um axioma de~$\Th{PA}$ e $= 0$ caso
contrário, podemos proceder assim: primeiro, verificamos se $!B$ é um
dos axiomas de~$\Th{Q}$. Se for, a resposta é ``sim'' e o valor de
$\Char{X}(!B) = 1$. Se não for, testamos se é uma instância do esquema
de indução. Isso pode ser feito sistematicamente; neste caso, talvez a
maneira mais fácil de perceber seja a seguinte: toda instância do
esquema de indução começa com certo número de quantificadores
universais e, depois, uma sub-!!{formula} que é um condicional. O
consequente desse condicional é
$\lforall[x][!A(x, y_1, \dots, y_n)]$, em que $x$ e $y_1$, \dots,
$y_n$ são todas as variáveis livres de~$!A$, e os quantificadores
iniciais de~$!B$ ligam as variáveis~$y_1$, \dots,~$y_n$. Depois de
extrair essa~$!A$ e verificar que suas variáveis livres correspondem às
variáveis ligadas pelos quantificadores universais iniciais e
por~$\lforall[x]$, passamos a verificar se o antecedente do condicional
corresponde a
\[
!A(\Obj 0, y_1, \dots, y_n) \land \lforall[x][(!A(x, y_1, \dots, y_n)
\lif !A(x', y_1, \dots, y_n))]
\]
Novamente, se corresponder, $!B$ é uma instância do esquema de indução;
se não corresponder, $!B$ não é.
\end{ex}

Para responder a essa questão---e à questão mais geral de quais teorias
são completas ou decidíveis---será útil considerar também a definição
a seguir. Lembre-se de que um conjunto $X$ é !!{enumerable} sse é vazio
ou existe !!a{surjective} função~$f \colon \Nat \to X$. Tal função é
chamada de enumeração de~$X$.

\begin{defn}
Um conjunto $X$ é chamado de \emph{!!{computably enumerable}}
(abreviadamente, !!{c.e.}) sse é vazio ou possui uma enumeração
computável.
\end{defn}

Além da !!{axiomatizability}, outra condição imposta às teorias às quais
os teoremas da incompletude se aplicam será que elas sejam fortes o
suficiente para provar fatos básicos sobre funções computáveis e
relações !!{decidable}. Por ``fatos básicos'', entendemos !!{sentence}s
que expressam quais são os valores de funções computáveis para cada um
de seus argumentos. E, por ``fortes o suficiente'', entendemos que as
teorias em questão contam essas sentenças entre seus teoremas.
Considere, por exemplo, uma função computável prototípica: a adição. O
valor de $+$ para os argumentos $2$ e $3$ é~$5$, isto é, $2+3 = 5$.
Uma sentença na linguagem da aritmética que expressa que o valor de $+$
para os argumentos $2$ e $3$ é~$5$ é:
$(\num{2} + \num{3}) = \num{5}$. E, por exemplo, $\Th{Q}$ prova essa
sentença. Mais geralmente, gostaríamos que, para cada função computável
$f(x_1, x_2)$, houvesse !!a{formula}~$!A_f(x_1, x_2, y)$ em
$\Lang{L_A}$ tal que $\Th{Q} \Proves
!A_f(\num{n_1}, \num{n_2}, \num{m})$ sempre que $f(n_1, n_2) = m$.
Desse modo, $\Th{Q}$ prova que o valor de~$f$ para os argumentos $n_1$,
$n_2$ é~$m$. Na verdade, exigimos que prove um pouco mais: que nenhum
outro número é o valor de~$f$ para os argumentos $n_1$,~$n_2$. E o
mesmo vale para relações !!{decidable}. Isso é formulado com precisão
nas duas definições a seguir.

\begin{defn}
Uma !!{formula}~$!A(x_1, \dots, x_k, y)$ \emph{!!{represents}} a
função $f\colon \Nat^k \to \Nat$ em~$\Gamma$ sse, sempre que
$f(n_1, \dots, n_k) = m$, então
\begin{enumerate}
\item $\Gamma \Proves !A(\num{n_1}, \dots, \num{n_k}, \num{m})$, e
\item $\Gamma \Proves \lforall[y](!A(\num{n_1}, \dots, \num{n_k},
y) \lif y = \num{m})$.
\end{enumerate}
\end{defn}

\begin{defn}
Uma !!{formula}~$!A(x_1, \dots, x_k)$ \emph{!!{represents}} a relação
$R \subseteq \Nat^k$ sse,
\begin{enumerate}
\item sempre que $R(n_1, \dots, n_k)$, $\Gamma \Proves
!A(\num{n_1}, \dots, \num{n_k})$, e
\item sempre que não $R(n_1, \dots, n_k)$, $\Gamma \Proves \lnot
!A(\num{n_1}, \dots, \num{n_k})$.
\end{enumerate}
\end{defn}

Uma teoria é ``forte o suficiente'' para que os teoremas da
incompletude se apliquem se ela !!{represents} todas as funções
computáveis e todas as relações !!{decidable}. $\Th{Q}$ e suas
extensões satisfazem essa condição, mas levaremos algum tempo para
estabelecê-la---trata-se de um fato não trivial sobre os tipos de coisa
que $\Th{Q}$ pode provar, e é difícil demonstrá-lo porque $\Th{Q}$ tem
apenas alguns axiomas a partir dos quais teremos de provar todos esses
fatos. Contudo, $\Th{Q}$ é uma teoria muito fraca. Assim, embora seja
difícil provar que $\Th{Q}$ representa todas as funções computáveis, a
maioria das teorias interessantes é mais forte que~$\Th{Q}$, isto é,
prova mais do que $\Th{Q}$. E, se $\Th{Q}$ prova alguma coisa, toda
teoria mais forte também a prova; como $\Th{Q}$ representa todas as
funções computáveis, toda teoria mais forte também as representa. Isso
significa que muitas teorias interessantes satisfazem essa condição
dos teoremas da incompletude. Portanto, nosso trabalho árduo dará
resultado, pois mostrará que os teoremas da incompletude se aplicam a
uma ampla classe de teorias. Certamente, toda teoria que pretenda
formalizar ``toda a matemática'' deve provar tudo o que $\Th{Q}$ prova,
pois deve, no mínimo, ser capaz de captar os resultados de computações
elementares. Assim, toda candidata a teoria de ``toda a matemática''
será uma teoria à qual se aplicam os teoremas da incompletude.


\end{document}
